% Quelle: 6. Decodierung von linearen Codes.pdf
% Art: Vorlesung
% Themen: Standardarray-Decodierung, Nebenklassen (Cosets), Klassenführer, Syndrom, Syndromdecodierung, Syndromdecodierung bei zyklischen Codes, Syndrompolynom, binärer Hamming-Code H2(3), modifizierter ISBN-Code, Fehlerstelle und Fehlergröße
\section{Vorlesung 6: Dekodierung von linearen Codes}
\subsection{6.1 Standardarray Dekodierung, Slepian (1960)}

Sei $C$ ein linearer, $e$ Fehler korrigierender Code $[n,k,d]_p$, $|C| = p^k$, $2e+1 \le d$.

\begin{tabular}{ll}
Gegeben: & empfangenes Wort $\mathbf{y} \in F_p^{\,n}$.\\
Gesucht: & Codewort $\mathbf{c} \in C$ mit
\end{tabular}
\[
  d(\mathbf{y},\mathbf{c}) = \min\{\, d(\mathbf{y},\mathbf{w}) \mid \mathbf{w} \in C \,\}
\]

\begin{satz}
Falls $\min\{\, d(\mathbf{y},\mathbf{w}) \mid \mathbf{w} \in C \,\} \le e$, dann ist $\mathbf{c}$ eindeutig.
\end{satz}
\begin{proof}
klar
\end{proof}

\begin{definition}[Nebenklasse (Coset)]
\[
  \mathbf{a} + C = \{\, \mathbf{a} + \mathbf{c} \mid \mathbf{c} \in C \,\}
\]
\end{definition}

\begin{satz}
\begin{enumerate}[label=(\arabic*)]
  \item Jedes Wort $\mathbf{y} \in F_p^{\,n}$ liegt in genau einer Nebenklasse,
  d.h.\ entweder die Nebenklassen $\mathbf{a}+C$, $\mathbf{y}+C$ stimmen überein oder sind disjunkt, also
  \[
    \text{entweder}\quad \mathbf{a}+C = \mathbf{y}+C \quad\text{oder}\quad \mathbf{a}+C \cap \mathbf{y}+C = \emptyset .
  \]
  \item $|\mathbf{a}+C| = |C| = p^k$
  \item Sei $s$ = Anzahl der Nebenklassen\\
  $\Rightarrow\ s \cdot |C| = p^n \ \Rightarrow\ s = p^{n-k}$
\end{enumerate}
\end{satz}

\begin{beispiel}
$C = \{\, 0000, 1011, 0101, 1110 \,\}$, $p = 2$.\\
$C$ ist linear, Parameter $[4,2,2]_2$, $d_C = 2$,
$G = \begin{pmatrix} 1 & 0 & 1 & 1 \\ 0 & 1 & 0 & 1 \end{pmatrix}$, \quad $|F_2^{\,4}| = 2^4 = 16$.

Die Nebenklassen sind
\begin{align*}
  0000 + C &= \{\, 0000, 1011, 0101, 1110 \,\}\\
  1000 + C &= \{\, 1000, 0011, 1101, 0110 \,\}\\
  0100 + C &= \{\, 0100, 1111, 0001, 1010 \,\}\\
  0010 + C &= \{\, 0010, 1001, 0111, 1100 \,\}
\end{align*}
\end{beispiel}

\begin{definition}[Klassenführer, Standardarray]
In jeder Nebenklasse gibt es ein Element $\mathbf{a}_i$ mit kleinstem Gewicht.
$\mathbf{a}_i$ heißt \textbf{Klassenführer}.

Obiges Schema von $F_p^{\,n}$ mit den Klassenführern in der ersten Spalte heißt

\begin{center}
\begin{tabular}{llllll}
\multicolumn{4}{l}{\textbf{Standardarray}} & & \textbf{Syndrom}\\
$\mathbf{o}$ & $\mathbf{c}_2$ & $\dots$ & $\mathbf{c}_t$ & & $\mathbf{o}$\\
$\mathbf{a}_2$ & $\mathbf{a}_2 + \mathbf{c}_2$ & $\dots$ & $\mathbf{a}_2 + \mathbf{c}_t$ & & $\mathrm{Syn}(\mathbf{a}_2)$\\
 & $\dots$ & & & & $\dots$\\
$\mathbf{a}_s$ & $\mathbf{a}_s + \mathbf{c}_2$ & $\dots$ & $\mathbf{a}_s + \mathbf{c}_t$ & & $\mathrm{Syn}(\mathbf{a}_s)$
\end{tabular}
\end{center}
mit $s = p^{n-k}$, $t = p^k$.
\end{definition}

\begin{algorithmus}[Standardarray Decodierung]
\begin{enumerate}[label=\arabic*. Schritt:, leftmargin=*]
  \item Suche das Wort $\mathbf{y} \in F_p^{\,n}$ im Standardarray.
  \[ \mathbf{y} = \mathbf{a}_i + \mathbf{c}_j \]
  \item Setze $\quad \mathbf{e} = \mathbf{a}_i,\quad \mathbf{c} = \mathbf{y} - \mathbf{e} = \mathbf{c}_j$
\end{enumerate}
Ergebnis: Falls $|\mathbf{a}_i| \le e$, ist $\mathbf{c}$ eindeutig.
\end{algorithmus}

\begin{beispiel}[Fortsetzung]
$C = \{\, 0000, 1011, 0101, 1110 \,\}$, $p = 2$.

\begin{center}
\begin{tabular}{lllll}
\toprule
Message & Codewort & empf.\ Wort & decodiertes Wort & decod.\ Message\\
\midrule
01 & $\mathbf{c} = 0101$ & $\mathbf{y} = 0001$ & $\mathbf{c}' = 0101$ & 01\\
01 & $\mathbf{c} = 0101$ & $\mathbf{y} = 0100$ & $\mathbf{c}'' = 0000$ & 00\\
\bottomrule
\end{tabular}
\end{center}
Solches passiert, wegen $d_C = 2$, also $d(\mathbf{y},\mathbf{c}') = d(\mathbf{y},\mathbf{c}'') = 1$.
\end{beispiel}

\begin{bemerkung}[Problem]
Das Standardarray ist sehr groß. $|F_p^{\,n}| = p^n$ Elemente.\\
Wie kann man den Aufwand für die Bestimmung der Nebenklasse von $\mathbf{y}$ minimieren?
\end{bemerkung}

\subsection{6.2 Syndromdecodierung}

\begin{definition}[Syndrom von $\mathbf{y}$]
\[ S(\mathbf{y}) = H\,\mathbf{y}^T \]
\end{definition}

\begin{satz}
Jedes Element derselben Nebenklasse besitzt dasselbe Syndrom:
\begin{enumerate}[label=(\arabic*)]
  \item $S(\mathbf{y}) = \mathbf{o} \quad\Leftrightarrow\quad \mathbf{y} \in C$
  \item $S(\mathbf{y}) = S(\mathbf{a}_i) \quad\Leftrightarrow\quad \mathbf{y} + C = \mathbf{a}_i + C$
\end{enumerate}
Das Syndrom $S(\mathbf{y})$ sagt uns in welcher Nebenklasse $\mathbf{y}$ zu finden ist.
Also brauchen wir vom Standardarray nur die 1.\ Spalte und eine zusätzliche letzte Spalte mit den Syndromen.
\end{satz}

\begin{algorithmus}[Syndromdecodierung]
\begin{enumerate}[label=\arabic*. Schritt:, leftmargin=*]
  \item Berechne das Syndrom $S(\mathbf{y})$.
  \item Suche das Syndrom $S(\mathbf{y}) = S(\mathbf{a}_i)$ in der Syndromspalte.
  \item Setze $\quad \mathbf{e} = \mathbf{a}_i,\quad \mathbf{c} = \mathbf{y} - \mathbf{e} = \mathbf{c}_j$
\end{enumerate}
Ergebnis: Falls $|\mathbf{a}_i| \le e$, ist $\mathbf{c}$ eindeutig.
\end{algorithmus}

\begin{bemerkung}
Falls $s = p^{n-k}$ zu groß ist, speichert man nur Nebenklassen mit $|\mathbf{a}_i| \le e$ und fordert ein erneutes Senden von $\mathbf{c}$, falls $\mathbf{y}$ nicht decodiert werden kann.
\end{bemerkung}

\subsubsection*{Syndromdecodierung bei zyklischen Codes}

\begin{bemerkung}[Syndromdecodierung bei zyklischen Codes]
Falls $C$ ein zyklischer Code mit Generatorpolynom $g(X)$, $\operatorname{grad} g(X) = r$, ist,
berechnet man das Syndrom von $y$ nicht über $S(y) = H\,\mathbf{y}^T$,
sondern man dividiert $y(X)$ mit Rest durch $g(X)$
\[
  y(X) = q(X)\,g(X) + s(X) \quad\text{mit}\quad \operatorname{grad} s(X) < r
\]
$s(X)$ ist das Syndrompolynom.\\
Die Bestimmung von $G$, $H$ und $h(X)$ kann entfallen.

Die binären Hamming Codes $H_2(m)$ sind äquivalent zu einem zyklischen Code.
\end{bemerkung}

\begin{beispiel}
$C$ = binärer Hammingcode $H_2(3)$ mit Parametern $[7,4,3]_2$, $g(X) = X^3 + X + 1$.\\
$\omega = X$ ist Nullstelle des primitiven Polynom $X^3 + X + 1 \in F_2[X]$.\\
Die Vektordarstellung von $X$ und ihren Potenzen liefern einen zyklischen Code
\[
  H = (\,1 \ \ X \ \ X^2 \ \ \dots \ \ X^{n-1}\,).
\]
\[
  X = \begin{pmatrix} 0\\1\\0 \end{pmatrix},\
  X^2 = \begin{pmatrix} 0\\0\\1 \end{pmatrix},\
  X^3 = X + 1 = \begin{pmatrix} 1\\1\\0 \end{pmatrix},\
  X^4 = X^2 + X = \begin{pmatrix} 0\\1\\1 \end{pmatrix},\
  X^5 = \begin{pmatrix} 1\\1\\1 \end{pmatrix},\
  X^6 = \begin{pmatrix} 1\\0\\1 \end{pmatrix}
\]
% Leere Einträge der Matrizen wie im Original (entsprechen 0).
\[
  H = \begin{pmatrix}
    1 &   &   & 1 & 0 & 1 & 1\\
      & 1 &   & 1 & 1 & 1 & 0\\
      &   & 1 & 0 & 1 & 1 & 1
  \end{pmatrix},\qquad
  G = \begin{pmatrix}
    1 & 1 &   & 1 &   &   &  \\
      & 1 & 1 &   & 1 &   &  \\
      &   & 1 & 1 &   & 1 &  \\
      &   &   & 1 & 1 &   & 1
  \end{pmatrix}
  = \begin{pmatrix} g_1\\ g_2\\ g_3\\ g_4 \end{pmatrix}
\]
Dann existieren $s = 2^{7-4} = 8$ Nebenklassen.

\begin{center}
\begin{tabular}{ll}
Klassenführer & Syndrom\\
$0$ & $0$\\
$1$ & $1$\\
$X$ & $X$\\
$X^2$ & $X^2$\\
$X^3$ & $X + 1$\\
$X^4$ & $X^2 + X$\\
$X^5$ & $X^2 + X + 1$\\
$X^6$ & $X^2 + 1$
\end{tabular}
\end{center}

Angenommen $y(X) = X^6 + X + 1$
\begin{align*}
  y(X) &= X^6 + X + 1 = (X^3 + X + 1)(X^3 + X + 1) + X^2 + X\\
  \Rightarrow\ s(X)) &= X^2 + X, \quad \text{Klassenführer } a_i(X) = X^4 % UNSICHER: doppelte Klammer "s(X))" so im Original
  \\
  \Rightarrow\ c(X) &= y(X) - a_i(X) = X^6 + X^4 + X + 1
\end{align*}
Tatsächlich ist
\[
  c(X) = X^6 + X^4 + X + 1 = (X^3 + 1)(X^3 + X + 1) = (X^3 + 1)\,g(X)
\]
und damit ein Codepolynom.\\
Geschrieben als Zeilenvektor
\begin{align*}
  \mathbf{y} &= 1\ 1\ 0\ 0\ 0\ 0\ 1\\
  \mathbf{c} &= 1\ 1\ 0\ 0\ 1\ 0\ 1 \qquad d(\mathbf{y},\mathbf{c}) = 1
\end{align*}
\end{beispiel}

\subsection{6.3 Decodierung eines modifizierten ISBN-Codes}

Der ISBN-Code ist ein modulo-11 Code $\subseteq F_{11}^{\,10}$. Seine Kontrollmatrix besteht nur aus der Zeile
$\mathbf{h} = (1\ 2\ 3\ \dots\ 10)$, die letzte Ziffer $c_{10}$ eines Codewortes $\mathbf{c}$ ist die Prüfziffer.
Um einen Fehler zu korrigieren, muss die Fehlerstelle $x$, $1 \le x \le 10$, und die Fehlergröße $e_x$ berechnet werden.
Dazu sind zwei Gleichungen bzw.\ zwei Prüfziffern erforderlich. Setzen
\[
  H = \begin{pmatrix} 1 & 1 & 1 & \dots & 1\\ 1 & 2 & 3 & \dots & 10 \end{pmatrix},
\]
Formen $H$ in systematische Form $(A \mid E)$ um (Einheitsmatrix rechts), um die Generatormatrix zu bestimmen.
\[
  \to\qquad (A \mid E) = \begin{pmatrix}
    9 & 8 & 7 & 6 & 5 & 4 & 3 & 2 & 1 & 0\\
    3 & 4 & 5 & 6 & 7 & 8 & 9 & 10 & 0 & 1
  \end{pmatrix},
\]
% Leere Einträge der Matrix wie im Original (entsprechen 0).
\[
  \Rightarrow\qquad G = (E \mid -A^T) = \begin{pmatrix}
    1 & & & & & & & & 2 & 8\\
    & 1 & & & & & & & 3 & 7\\
    & & 1 & & & & & & 4 & 6\\
    & & & 1 & & & & & 5 & 5\\
    & & & & 1 & & & & 6 & 4\\
    & & & & & 1 & & & 7 & 3\\
    & & & & & & 1 & & 8 & 2\\
    & & & & & & & 1 & 9 & 1
  \end{pmatrix}
\]
\begin{align*}
  \Rightarrow\qquad & \mathbf{c} = \mathbf{a}\,G, \quad \text{also } c_i = a_i \text{ für } 1 \le i \le 8.\\
  & c_9 \equiv \sum_{i=1}^{8} (1+i)\,c_i \pmod{11}, \qquad c_{10} \equiv \sum_{i=1}^{8} (9-i)\,c_i \pmod{11}.
\end{align*}

Sei $\mathbf{r}$ wie in Abschnitt 1.3 das empfangene Wort. Falls $\mathbf{r} = \mathbf{c}$, dann gilt
\[
  \mathbf{s}(\mathbf{r}) = H\,\mathbf{r}^T = H\,\mathbf{c}^T = \mathbf{o}
\]
Falls ein Fehler an der Fehlerstelle $x$ mit der Größe $e_x$ auftrat, sieht $\mathbf{r}$ so aus
\begin{align*}
  & \mathbf{r} = \mathbf{c} + \mathbf{e} \quad\text{mit}\quad \mathbf{e} = (0\ \dots\ 0\ e_x\ 0\ \dots\ 0)\\
  \Rightarrow\qquad & H\,\mathbf{r}^T = H\,\mathbf{e}^T = \begin{pmatrix} e_x\\ x e_x \end{pmatrix}
    = \begin{pmatrix} 1\\ x \end{pmatrix} e_x = \begin{pmatrix} s_1\\ s_2 \end{pmatrix}\\
  & s_1 = e_x,\quad s_2 = x\,e_x,\quad x = \frac{s_2}{s_1},\quad c_x = r_x - e_x
\end{align*}
d.h.\ $s_1$ ist die Fehlergröße und die Fehlerstelle ist $x = s_2 / s_1$.

\begin{beispiel}
Sei $\mathbf{a} = 0\ 0\ \dots\ 0$, 8 Nullen, dann ist $\mathbf{c} = 0\ 0\ \dots\ 0\ 0\ 0$, 10 Nullen.\\
Angenommen
\[ \mathbf{r} = 0\ 0\ \dots\ 0\ 5\ 0. \]
Dann ist $S(\mathbf{r}) = H\,\mathbf{r}^T = \begin{pmatrix} 5\\ 45 \end{pmatrix} = \begin{pmatrix} 1\\ 9 \end{pmatrix} 5$, also $e_x = 5$, $x = 9$.\\
Tatsächlich ist $\mathbf{c} = \mathbf{o}$ Codewort in jedem linearen Code.
\end{beispiel}

\begin{beispiel}
Sei $\mathbf{r} = 3\ 5\ \dots\ 0\ 0\ 0$.\\
Dann ist $\mathbf{s} = H\,\mathbf{r}^T = \begin{pmatrix} 8\\ 13 \end{pmatrix} = \begin{pmatrix} 8\\ 2 \end{pmatrix}$, also $e_x = 8$, $x = 2/8 = 1/4 = 3$.\\
Wegen $3 \cdot 4 = 12 = 1$, also $c_3 = 0 - 8 = 3$
\[ \mathbf{c} = 3\ 5\ 3\ \dots\ 0\ 0\ 0. \]
% Im Original als "1.3 + 2.5 + 3.3" (Punkt als Malzeichen) geschrieben.
Tatsächlich ist $H\,\mathbf{c}^T = \begin{pmatrix} 3+5+3\\ 1\cdot 3 + 2\cdot 5 + 3\cdot 3 \end{pmatrix} = \begin{pmatrix} 11\\ 22 \end{pmatrix} = \begin{pmatrix} 0\\ 0 \end{pmatrix}$, also $\mathbf{c}$ ein Codewort.
\end{beispiel}

\begin{bemerkung}
Falls $s_1 = 0$ und $s_2 \neq 0$, enthält $r$ mehr als einen Fehler.
\end{bemerkung}
